\section{Introduction}
\label{sec:intro}

Multi-agent systems powered by large language models (LLMs) are commonly
designed around explicit organizational structures and role
assignments~\cite{hong2024metagpt,qian2024chatdev,wu2024autogen}. While
formal diagrams suggest a clear division of labor among components such as
Planners, Executors and Reviewers, the actual execution can diverge from
these intentions~\cite{cemri2025why}. Because the roles are instantiated
only through natural language prompts, agents can ignore their constraints,
leading to a functional collapse in which every agent solves the whole
problem. An accuracy score cannot distinguish such a team from one that
works exactly as drawn. This disconnect raises a fundamental question: what
do assigned roles actually change in practice?

In our experiments, roles do produce the division of labor they prescribe,
yet it buys no better first answer. If specialization is not where roles
pay off, their effect must lie elsewhere. We argue that roles also allocate
decision rights and thereby shape the flow of external information. In a
Planner-Executor-Reviewer (\PER{}) team, the Reviewer acts as a gatekeeper:
by approving a candidate or demanding further deliberation, it decides when
the team consults an external evaluator. This decision is crucial because
the evaluator returns a verdict (PASS or FAIL) and,
after a FAIL, a repair hint, corrective information that internal
deliberation cannot generate. A cautious Reviewer saves evaluations, but under a capped budget
every withheld submission is a missed opportunity to recover from an
error.

To isolate these effects, we hold the actor model (the language model that
plays all three agents), the agent topology, the agent-call budget and the
ground-truth evaluator constant, and vary the organization in two
directions: we ablate every role instruction, and we enforce the division
of labor through typed message formats. Because feedback reaches each agent
through its prompt, a role ablation also rewrites the feedback channel, so
we document the information each condition receives (\S\ref{sec:method}).
In our main comparison both teams have equal access to the evaluator's
verdict and hint, but the realized number and content of verdicts differ,
and this difference is what roles control. Our findings are as follows.
First, roles induce functional specialization, but it does not improve
first-pass accuracy, consistent with prior role
ablations~\cite{dylan2026equal,orogat2026architectural}. Second, roles
create a gatekeeping effect. A role-structured Reviewer often withholds
approval, frequently before any feedback exists, so its team submits less
and receives fewer corrective verdicts; in an exploratory comparison with
equal access to feedback, role-free teams recover from a failed first answer
more often in all five model and benchmark settings and reach higher final
accuracy. Third, the gap stems mainly from feedback opportunities rather
than from how well each team uses feedback: it vanishes in a post hoc
comparison at equal numbers of verdicts, and, in a single exploratory run, forcing the
role-structured team to submit in every cycle (one pass of the three
agents) raises its recovery to the role-free level for the gemma-4 model. Fourth, two methodological
cautions follow: a role ablation that silently removes the feedback channel
can make roles look necessary for repair, and enforcing roles through typed
interfaces loosens the gate but leaves $8\%$ to $17\%$ of enforced trajectories unfinished, so accuracy does not rise. Finally, evaluating multi-agent
organizations requires reporting recovery after a failed first attempt and
treating decision rights and information pathways as part of the
organization's specification (\S\ref{sec:discussion}).
